\documentclass[12pt,a4paper]{article}

\usepackage[utf8]{inputenc}
\usepackage[T1]{fontenc}
\usepackage[polish]{babel}
\usepackage{polski}
\usepackage{geometry}
\geometry{margin=2.5cm}
\usepackage{graphicx}
\usepackage{hyperref}
\usepackage{enumitem}
\usepackage{listings}
\usepackage{booktabs}
\usepackage{float}

% ustawienia wygladu kodu
\lstset{
    basicstyle=\ttfamily\small,
    breaklines=true,
    frame=single,
    captionpos=b,
    extendedchars=true,
    keepspaces=true,
    columns=fixed,
    commentstyle=\normalfont\ttfamily,
    showstringspaces=false,
    % zeby polskie znaki w kodzie nie wywalaly bledow
    literate={ą}{{\k{a}}}1 {ć}{{\'c}}1 {ę}{{\k{e}}}1 {ł}{{\l{}}}1 {ń}{{\'n}}1 {ó}{{\'o}}1 {ś}{{\'s}}1 {ź}{{\'z}}1 {ż}{{\.z}}1 {Ą}{{\k{A}}}1 {Ć}{{\'C}}1 {Ę}{{\k{E}}}1 {Ł}{{\L{}}}1 {Ń}{{\'N}}1 {Ó}{{\'O}}1 {Ś}{{\'S}}1 {Ź}{{\'Z}}1 {Ż}{{\.Z}}1
}

% informacje o autorach
\title{Dokumentacja Funkcjonalna i Implementacyjna Projektu C\\ \large System wyznaczania współrzędnych węzłów grafów planarnych}
\author{Autorzy: Mateusz Bombola, Kajetan Karwacki}
\date{\today}

\begin{document}

\maketitle
\newpage
\tableofcontents
\newpage

% ========================================================
% CZĘŚĆ 1: DOKUMENTACJA FUNKCJONALNA
% ========================================================
\part{Dokumentacja Funkcjonalna}

\section{Cel projektu}
Celem projektu jest stworzenie aplikacji w języku C, która na podstawie topologii grafu planarnego wyznacza optymalne współrzędne dwuwymiarowe $(x, y)$ dla jego wierzchołków. Program służy jako silnik obliczeniowy przetwarzający dane o połączeniach w celu wygenerowania układu umożliwiającego czytelną wizualizację grafu.

\section{Opis algorytmów rozmieszczania grafu}
W programie zaimplementowano dwie niezależne metody wyznaczania współrzędnych:

\begin{itemize}
    \item \textbf{Algorytm SMACOF (\texttt{smacof}):} Służy do rozmieszczania elementów w przestrzeni na podstawie ich wzajemnych podobieństw lub odległości. Działa poprzez iteracyjne poprawianie położenia punktów tak, aby odległości między nimi jak najlepiej odpowiadały zadanym wartościom. W efekcie powstaje układ, w którym elementy bardziej podobne znajdują się bliżej siebie, a mniej podobne dalej.
    
    \item \textbf{Algorytm Fruchtermana-Reingolda (\texttt{fr}):} Automatycznie rozmieszcza elementy grafu w przestrzeni tak, aby był on czytelny i przejrzysty. Traktuje węzły jak obiekty, które wzajemnie się odpychają, a połączone elementy przyciągają się jak sprężyny. Dzięki wielokrotnemu przeliczaniu tych zależności program stopniowo ustala układ, w którym struktura grafu jest łatwa do zrozumienia.
\end{itemize}

\section{Interfejs linii poleceń}
Program jest sterowany poprzez argumenty przekazywane w wierszu poleceń. Składnia wywołania:
\begin{center}
    \begin{verbatim}
    ./graph_layout -i <input> -o <output> -a <smacof|fr> [-n iter] [-f txt|bin]
    \end{verbatim}
\end{center}

\subsection{Lista dostępnych flag}
\begin{table}[H]
\centering
\begin{tabular}{@{}lll@{}}
\toprule
\textbf{Flaga} & \textbf{Argument} & \textbf{Opis} \\ \midrule
\texttt{-i} & \textit{ścieżka} & Ścieżka do wejściowego pliku tekstowego (wymagane). \\
\texttt{-o} & \textit{ścieżka} & Ścieżka do pliku wynikowego (wymagane). \\
\texttt{-a} & \textit{algorytm} & Wybór algorytmu: \texttt{smacof} lub \texttt{fr}. \\
\texttt{-n} & \textit{liczba} & Liczba iteracji (domyślnie 1000). \\
\texttt{-f} & \textit{format} & Format wyjściowy: \texttt{txt} lub \texttt{bin}. \\
\texttt{-h} & --- & Wyświetlenie pomocy technicznej. \\
\bottomrule
\end{tabular}
\caption{Parametry sterujące programem}
\end{table}

\section{Formaty danych}
\subsection{Format wejściowy}
Każda linia reprezentuje krawędź: \texttt{<nazwa> <od\_id> <do\_id> <waga>}.
\\
\textbf{Przykład:}
\begin{lstlisting}
AB   1  2  1
BC   2  3  1
CD   3  4  1
DB   4  2  1.407
\end{lstlisting}

\subsection{Format wyjściowy}
Wynikiem działania programu ma być plik tekstowy lub binarny (w zależności od wyboru użytkownika) zawierający listę współrzędnych wierzchołków: \texttt{<wierzchołek> <współrzędna\_x> <współrzędna\_y>}.
\\
\textbf{Przykład:}
\begin{lstlisting}
1 0.0 0.0
2 1.0 0.0
3 1.0 1.0
4 0.0 1.0
\end{lstlisting}

\section{Kody błędów}
\begin{itemize}
    \item \textbf{Kod 1:} Błąd otwarcia pliku wejściowego.
    \item \textbf{Kod 2:} Nieprawidłowy format danych wejściowych.
    \item \textbf{Kod 3:} Nieznany algorytm (oczekiwano \texttt{smacof} lub \texttt{fr}).
    \item \textbf{Kod 4:} Błąd alokacji pamięci.
\end{itemize}

\section{Przykłady użycia}
% gotowe komendy do skopiowania

1. Wyznaczenie współrzędnych metodą Fruchtermana-Reingolda i zapis do tekstu:
\begin{lstlisting}[language=bash]
./graph_layout -i dane.txt -o wynik.txt -a fr -n 1000 -f txt
\end{lstlisting}

2. Wyznaczenie współrzędnych metodą SMACOF i zapis binarny (500 iteracji):
\begin{lstlisting}[language=bash]
./graph_layout -i siec.txt -o mapa.bin -a smacof -n 500 -f bin
\end{lstlisting}

\newpage
% ========================================================
% CZĘŚĆ 2: DOKUMENTACJA IMPLEMENTACYJNA
% ========================================================
\part{Dokumentacja Implementacyjna}

\section{Architektura systemu i podział na moduły}
System został zaprojektowany z zachowaniem zasad programowania modułowego. Kod źródłowy w języku C rozdzielono na osobne pliki, oddzielając logikę sterującą od implementacji poszczególnych algorytmów obliczeniowych. Dodatkowo zintegrowano środowisko testowe oparte na skryptach Python.

\subsection{Moduły języka C}
\begin{itemize}
    \item \textbf{\texttt{graph.h}:} Główny plik nagłówkowy spinający projekt. Zawiera definicje struktur danych (\texttt{Vertex}, \texttt{Edge}), stałych matematycznych (np. \texttt{EPSILON}) oraz deklaracje prototypów funkcji obliczeniowych.
    
    \item \textbf{\texttt{main.c}:} Moduł inicjujący i sterujący. Parsuje argumenty linii poleceń , wczytuje i weryfikuje dane topologii, dynamicznie alokuje pamięć, inicjalizuje wstępne współrzędne, a po wykonaniu obliczeń zapisuje wyniki do plików.
    
    \item \textbf{\texttt{smacof.c}:} Niezależny moduł realizujący algorytm SMACOF. Główną logikę umieszczono w funkcji \texttt{calculate\_smacof}, wspieranej przez matematyczną funkcję pomocniczą \texttt{safe\_dist}.
    
    \item \textbf{\texttt{fruchterman.c}:} Moduł zawierający implementację siłowego algorytmu optymalizacyjnego wewnątrz funkcji \texttt{calculate\_fruchterman}.
\end{itemize}

\subsection{Moduły testowe i narzędziowe}
Na potrzeby zapewnienia jakości aplikacji przygotowano zestaw narzędzi:
\begin{itemize}
    \item \textbf{\texttt{generator.py}:} Generator pseudolosowych instancji grafów. Generuje siatki połączeń dla zadanej liczby węzłów i krawędzi (z opcją losowania wag), produkując pliki wejściowe w formacie akceptowalnym przez program.
    \item \textbf{\texttt{wizualizacja.py}:} Analizator graficzny wykorzystujący bibliotekę \texttt{matplotlib}. Wczytuje oryginalną siatkę krawędzi oraz wygenerowane pozycje X i Y, umożliwiając wizualną weryfikację węzłów na płaszczyźnie.
\end{itemize}

\section{Struktury danych}
Podstawowymi strukturami wykorzystywanymi do reprezentacji grafów są:

\begin{lstlisting}[language=C]
typedef struct { 
    int id; 
    double x, y; 
} Vertex;

typedef struct { 
    int u, v; 
    double weight; 
} Edge;
\end{lstlisting}

\textbf{Mechanizm mapowania węzłów (\texttt{main.c}):} Numery węzłów (identyfikatory) podane w pliku wejściowym nie zawsze występują po kolei (np. mogą to być węzły 10 i 100 bez żadnych węzłów pomiędzy). Moduł parsujący buduje strukturę \texttt{node\_map}, przypisując takim węzłom kolejne, ciągłe indeksy ($[0, N-1]$). Pozwala to programowi optymalnie używać pamięci i zapobiega tworzeniu ogromnych, pustych tablic. Dopiero przy zapisie wyników do pliku, program przywraca węzłom ich oryginalne identyfikatory.

\section{Implementacja algorytmów optymalizacyjnych}

\subsection{SMACOF (\texttt{smacof.c})}
Zastosowano algorytm oparty na transformacji Guttmana, składający się z kluczowych faz:
\begin{itemize}
    \item \textbf{Macierz odległości:} Na bazie wczytanych krawędzi algorytm najpierw używa metody \textbf{Floyda-Warshalla} do propagacji najkrótszych ścieżek między wszystkimi niesąsiadującymi węzłami grafu, obsługując izolowane klastry dużą stałą odległością.
    \item \textbf{Inicjalizacja współrzędnych:} Zamiast losować pozycje startowe, program na samym początku układa wszystkie węzły równomiernie na obwodzie koła (korzystając z funkcji \texttt{sin()} i \texttt{cos()}). Zapobiega to nakładaniu się punktów na siebie i pomaga algorytmowi płynniej znaleźć optymalny układ, bez ryzyka utknięcia w złym ułożeniu początkowym.
    \item \textbf{Zabezpieczenie przed dzieleniem przez zero:} Jeśli dwa węzły znajdą się dokładnie w tym samym punkcie (odległość wyniesie 0), wzory algorytmu wygenerowałyby błąd matematyczny psujący wyniki. Aby temu zapobiec, wprowadzono funkcję \texttt{safe\_dist}, która podmienia zerową odległość na minimalną dopuszczalną wartość (\texttt{EPSILON = 1e-9}), zapewniając stabilność całego programu.
\end{itemize}

\subsection{Fruchterman-Reingold (\texttt{fruchterman.c})}
Zaimplementowano klasyczny model sprężyn i odpychających się ładunków elektrycznych. Obliczenia na każdą iterację \texttt{iterations} dzielą się na:
\begin{itemize}
    \item Faza $O(V^2)$: globalne odpychanie poszczególnych wierzchołków na zasadzie odwrotnej proporcjonalności do ich odległości.
    \item Faza $O(E)$: iteracja wzdłuż bezpośrednich krawędzi (zdefiniowanych w wektorze struktur \texttt{Edge}) i przyciąganie połączonych wierzchołków wzajemnymi wektorami sił z jednoczesnym uwzględnieniem wagi \texttt{(weight)}.
\end{itemize}

\section{Specyfikacja wejścia i wyjścia}
Poniższa specyfikacja definiuje formaty danych wymienianych między modułem obliczeniowym, a zewnętrznymi komponentami systemu.

\subsection{Wejście}
Program parsuje pliki wejściowe liniowo wywołując instrukcję \texttt{fscanf}. Kod zawarty w \texttt{main.c} aplikuje automatyczną naprawę błędnych danych. Dowolne wejściowe wagi krawędzi $\le 0$ są interpretowane i nadpisywane systemową stałą o wartości \texttt{0.1}, co zapobiega propagacji błędu przez ujemne odległości.

\subsection{Format plików tekstowych (\texttt{txt})}
Wynikowa lokalizacja wierzchołków generowana jest do postaci ASCII i powinna być parsowana wyrażeniem: \verb|%d %f %f\n| (Kolejno: Integer ID, Double Oś X, Double Oś Y).

\subsection{Format plików binarnych (\texttt{bin})}
Opcjonalny zrzut strukturalny oparty o procedurę \texttt{fwrite} (\texttt{save\_to\_binary}). Wymaga zastosowania odpowiednich technik deserializacji:
\begin{enumerate}
    \item \textbf{Liczba wierzchołków ($N$):} Pierwsze 4 bajty pliku to zwykła liczba całkowita (typ \texttt{int}). Informuje ona, ile łącznie węzłów zapisano w pliku.
    \item \textbf{Dane pojedynczego wierzchołka:} Reszta pliku natomiast wypełniona jest blokami 20-bajtowymi. Architektura bloku to: $4$ bajty przestrzeni \texttt{id} (Signed Integer 32-bit), $8$ bajtów \texttt{x} (Double 64-bit) oraz $8$ bajtów \texttt{y} (Double 64-bit).
\end{enumerate}

\section{Zarządzanie pamięcią i łączenie}
Kluczowe wektory i macierze przetrzymywane są na stercie. Główna tablica struktur węzłowych \texttt{vertices} jest alokowana z użyciem \texttt{malloc} przez \texttt{main.c} (tylko na podstawie wyliczonej ilości realnie zidentyfikowanych wierzchołków, nie na stałym przedziale rozmiaru grafu). Ponadto tablice dwuwymiarowe w funkcji \texttt{calculate\_smacof} (odległości i wagi) są budowane jako macierze wskaźników dynamicznych.

Wszystkie dynamicznie zaalokowane zasoby są zwalniane za pomocą wywołań funkcji \texttt{free()} po zakończeaniu działania algorytmów lub tuż przed zamknięciem programu, co zapobiega wyciekom pamięci.

\textbf{Komenda kompilacji z wymuszeniem linkowania bibliotek współdzielonych dla trygonometrii:}
\begin{lstlisting}[language=bash]
gcc -o graph_layout main.c smacof.c fruchterman.c -lm
\end{lstlisting}

\newpage
% ========================================================
% CZĘŚĆ 3: DOKUMENTACJA PROJEKTOWA (KOŃCOWA)
% ========================================================
\part{Dokumentacja Projektowa}

\section{Wstęp i podsumowanie realizacji}
Niniejszy dokument stanowi zwieńczenie prac nad silnikiem obliczeniowym w języku C. Zgodnie z wymaganiami projektowymi, program został w pełni zaimplementowany, przetestowany i jest gotowy do integracji z docelowym interfejsem graficznym w języku Java. W ramach prac zrealizowano i przetestowano dwa niezależne algorytmy wyznaczania współrzędnych: model siłowy Fruchtermana-Reingolda oraz algorytm SMACOF oparty na skalowaniu wielowymiarowym.

\section{Metodologia testowania}
Aby rzetelnie porównać zaimplementowane algorytmy, wykorzystano dedykowane środowisko testowe oparte na skryptach w języku Python:
\begin{itemize}
    \item \textbf{Skrypt generujący (\texttt{generator.py}):} Pozwolił na masowe tworzenie grafów o różnej gęstości, liczbie węzłów oraz losowych wagach krawędzi. Umożliwiło to zbadanie stabilności programu w języku C przy dużych obciążeniach (powyżej 500 wierzchołków).
    \item \textbf{Skrypt wizualizacyjny (\texttt{wizualizacja.py}):} Moduł oparty na bibliotece \texttt{matplotlib}, wizualizujący wygenerowane przez program C współrzędne na płaszczyźnie, co pozwoliło na obiektywną ocenę jakości estetycznej układu grafu.
\end{itemize}

\section{Porównanie algorytmów i wyniki testów}

\subsection{Analiza wydajności i czasu wykonania}
Zarówno SMACOF, jak i algorytm Fruchtermana-Reingolda wykazują odmienną złożoność obliczeniową. Poniższa tabela przedstawia uśrednione czasy wykonania dla różnych rozmiarów grafów (przy stałej liczbie 1000 iteracji):

\begin{table}[H]
\centering
\begin{tabular}{@{}lccc@{}}
\toprule
\textbf{Algorytm} & \textbf{Mały graf (20 węzłów)} & \textbf{Średni graf (100 węzłów)} & \textbf{Duży graf (500 węzłów)} \\ \midrule
Fruchterman-Reingold & <UZUPEŁNIJ> ms & <UZUPEŁNIJ> ms & <UZUPEŁNIJ> ms \\
SMACOF & <UZUPEŁNIJ> ms & <UZUPEŁNIJ> ms & <UZUPEŁNIJ> ms \\
\bottomrule
\end{tabular}
\caption{Porównanie wydajności algorytmów (1000 iteracji)}
\end{table}

\textbf{Wniosek wydajnościowy:} 
<UZUPEŁNIJ: Napisz krótko, który algorytm był szybszy. Zazwyczaj FR jest szybszy, bo SMACOF musi na początku wyliczyć macierz Floyda-Warshalla dla wszystkich węzłów, co przy dużych grafach mocno obciąża procesor.>

\subsection{Ocena jakości wizualnej}
Poniższe grafiki przedstawiają ten sam graf wejściowy przetworzony przez oba zaimplementowane algorytmy.

% INSTRUKCJA: Zróbcie dwa screeny z Waszego programu w Pythonie (jeden z wynikiem FR, drugi SMACOF). 
% Zapiszcie je w folderze z tym plikiem LaTeX jako 'wynik_fr.png' i 'wynik_smacof.png'.
% Następnie odkomentujcie (usuńcie znaki %) z poniższego bloku kodu:

% \begin{figure}[H]
%     \centering
%     \begin{minipage}{0.45\textwidth}
%         \centering
%         \includegraphics[width=\linewidth]{wynik_fr.png}
%         \caption{Wynik alg. Fruchtermana-Reingolda}
%     \end{minipage}\hfill
%     \begin{minipage}{0.45\textwidth}
%         \centering
%         \includegraphics[width=\linewidth]{wynik_smacof.png}
%         \caption{Wynik alg. SMACOF}
%     \end{minipage}
% \end{figure}

Wizualna analiza wygenerowanych układów wykazuje, że:
\begin{itemize}
    \item \textbf{Fruchterman-Reingold} lepiej radzi sobie z równomiernym rozkładaniem wierzchołków, tworząc symetryczne układy ułatwiające analizę topologii.
    \item \textbf{SMACOF} wykazuje wyższą precyzję w odzwierciedlaniu różnic w wagach krawędzi (zachowanie proporcji metrycznych), jednak przy grafach bardzo gęstych ma tendencję do skupiania wielu węzłów w ciasne klastry.
\end{itemize}

\section{Napotkane problemy i zastosowane rozwiązania}
Podczas implementacji zidentyfikowano i rozwiązano następujące wyzwania techniczne:
\begin{enumerate}
    \item \textbf{Niestabilność numeryczna przy nakładających się wierzchołkach:} W obu modelach nakładające się na siebie węzły (odległość równa 0) powodowały dzielenie przez zero. Rozwiązaniem było zaimplementowanie funkcji \texttt{safe\_dist}, podmieniającej zerową odległość na minimalną stałą \texttt{EPSILON}.
    \item \textbf{Optymalizacja zużycia pamięci:} Pliki testowe posiadały rzadkie, nieciągłe identyfikatory (np. węzeł 10 i 1000). Bezpośrednia alokacja doprowadziłaby do ogromnego marnotrawstwa pamięci. Zastosowano mechanizm wektorowego mapowania (\texttt{node\_map}), który spłaszczył indeksy na ciągły obszar roboczy.
    \item \textbf{Problem minimów lokalnych:} Losowa inicjalizacja współrzędnych często powodowała zatrzymanie się algorytmów w nieestetycznym, zaplątanym układzie. Zmiana geometrii startowej na równomierne rozmieszczenie wierzchołków na okręgu znacząco zwiększyła odsetek udanych optymalizacji.
\end{enumerate}

\section{Rekomendacje dla zespołu JAVA}
Na podstawie wyników testów, zespołowi odpowiedzialnemu za implementację interfejsu graficznego (GUI) rekomendujemy:
\begin{itemize}
    \item Ustawienie algorytmu \textbf{Fruchtermana-Reingolda jako domyślnego} narzędzia do wizualizacji, ze względu na generowanie czytelniejszych dla ludzkiego oka układów symetrycznych.
    \item Umożliwienie użytkownikowi dynamicznej zmiany liczby iteracji z poziomu panelu bocznego. Dla bardzo dużych sieci, redukcja iteracji z 1000 na 300 pozwoli zachować płynność działania aplikacji.
    \item Preferowanie binarnego formatu wymiany danych (\texttt{.bin}) w przypadku ładowania grafów zawierających powyżej 1000 połączeń, co pozwoli znacząco skrócić czas ładowania (deserializacji).
\end{itemize}

\end{document}